\addtocontents{toc}{\protect\setcounter{tocdepth}{0}}
\chapter[Anexo: Resultados de Iteraciones]{ANEXO: Resultados de Iteraciones en el Desarrollo}

El desarrollo se organiza en tres sprints. Cada uno termina con una versión
instalable que se revisa con la contraparte de Enfermería. Este anexo resume la
funcionalidad abordada en cada sprint y sus resultados.

\section{Sprint 1: flujo principal (26 de agosto al 15 de septiembre)}

\subsection*{Funcionalidad abordada}

Carga inicial y bienvenida (HU-01), creación de cuenta (HU-02), inicio de sesión
y recuperación de acceso (HU-03), cifrado de datos sensibles (HU-04), perfil del
cuidador (HU-05), gestión y selección de pacientes (HU-06 y HU-07), registro
clínico con clasificación de alertas (HU-09), historial (HU-10) y edición y
eliminación de registros (HU-11).

\subsection*{Resultados y observaciones}

Se completan las diez historias del sprint. El motor de reglas clasifica cada
registro según los umbrales de temperatura, frecuencia cardíaca, saturación de
oxígeno, frecuencia respiratoria y escala ESAS, y sus 40 casos de prueba cubren
los valores límite de cada umbral. El formulario incorpora además un tope
diario de registros programados, configurable por paciente, que distingue las
valoraciones rutinarias de las extra.

\section{Sprint 2: módulos funcionales (16 de septiembre al 20 de octubre)}

\subsection*{Funcionalidad abordada}

El sprint se organiza en cinco semanas, cada una en su propia rama del
repositorio.

\begin{table}[H]
\centering
\small
\tablazebra
\caption{Avance del sprint 2 por semana.}
\label{tab:sprint2-semanas}
\begin{tabularx}{\textwidth}{P{0.45}P{1.85}P{0.70}}
\toprule
\filaencabezado \encab{Semana} & \encab{Funcionalidad} & \encab{Historias} \\
\midrule
1 & Panel principal y exportación del historial a PDF y Excel. & HU-08, HU-12, HU-13 \\
2 & Biblioteca educativa con materiales descargables. & HU-19, HU-20 \\
3 & Videos, guías, infografías y contenido sin
conexión. & HU-21, HU-22 \\
4 & Asistente de orientación, preguntas frecuentes, conversaciones y favoritos
de la biblioteca. & HU-14, HU-15, HU-16 \\
5 & Recordatorios, notificaciones locales y conexión final de rutas. & HU-17, HU-18 \\
\bottomrule
\end{tabularx}
\end{table}

\subsection*{Resultados y observaciones}

El panel principal se aprueba en la revisión de la primera semana. El asistente
se implementa con emparejamiento por palabras clave sobre preguntas frecuentes
validadas, sin inteligencia artificial generativa, para asegurar que cada
respuesta provenga de contenido revisado por Enfermería. El servicio de
notificaciones queda implementado y probado con dobles de prueba, pero falta
validar en un dispositivo real la solicitud de permisos, la apertura desde la
notificación y la reprogramación al iniciar sesión.

\section{Sprint 3: integración (21 de octubre al 13 de noviembre)}

\subsection*{Funcionalidad abordada}

Separación de datos por paciente (HU-24), aislamiento de datos entre usuarios
(HU-25), navegación definitiva (HU-26) y funcionamiento sin conexión con
sincronización (HU-27).

\subsection*{Resultados y observaciones}

Las cuatro historias parten con una base funcional. Las consultas ya filtran
por paciente activo y el caché local de datos está activo. Falta
agregar la verificación en el servidor que formaliza el aislamiento, la cola de
sincronización de registros creados sin conexión y el aviso visual de estado de
conexión. La navegación actual usa una barra de cinco destinos, que el sprint
revisa frente al diseño propuesto de cuatro pestañas con un botón flotante.

\section{Migración al servidor propio} % REVISAR: fechas de la iteración

\subsection*{Funcionalidad abordada}

Reemplazo de la plataforma de servicios externa por un servidor propio: una API
REST (Node.js y Express) sobre PostgreSQL, ejecutada con Docker Compose y
destinada a contenedores Podman de la universidad. La aplicación pasa a usar un
cliente HTTP con sesión mediante JWT, la clave de datos se entrega desde el
servidor sellada con una clave maestra, los correos de recuperación y de
verificación se envían por SMTP y la biblioteca se sirve desde el propio
servidor. Fecha de la iteración: \TBD{fechas de la migración}.

\subsection*{Resultados y observaciones}

El servidor suma 54 pruebas automatizadas contra una base PostgreSQL real, y el
recorrido de punta a punta (registro, inicio de sesión, paciente y correos) se
probó en un teléfono Android. El cifrado en el teléfono se mantiene: el servidor
solo guarda texto cifrado y rechaza los campos clínicos en claro. El
funcionamiento sin conexión se rehízo con una cola persistente de escrituras
con reintentos, y la actualización en pantalla se obtiene consultando al servidor
cada 10 segundos y tras cada escritura. Queda pendiente el despliegue en el
servidor de la universidad, que depende de la activación de la cuenta.
